% ============================================================================
% CHAPTER 10: MODEL CLASSIFICATION SYSTEM
% ============================================================================

\chapter{Model Classification System}

\section{Classification System Architecture}

\subsection{3-Tier Ensemble Design}

The platform implements a sophisticated 3-tier cascade ensemble for optimal balance of speed, accuracy, and cost.

\begin{figure}[h]
\centering
\begin{tikzpicture}[
    node distance=1.5cm,
    tier/.style={rectangle, rounded corners, minimum width=4cm, minimum height=1.2cm, text centered, draw=black, fill=primaryblue!30, font=\bfseries},
    decision/.style={diamond, minimum width=2cm, minimum height=2cm, text centered, draw=black, fill=warningyellow!30, font=\small},
    result/.style={rectangle, rounded corners, minimum width=3cm, minimum height=1cm, text centered, draw=black, fill=successgreen!30}
]
    \node (input) [tier, fill=lightgray] {Document Input};
    \node (tier1) [tier, below of=input] {Tier 1: Rule-Based};
    \node (check1) [decision, below of=tier1] {Conf > 0.85?};
    \node (tier2) [tier, right of=check1, xshift=3.5cm] {Tier 2: BERT};
    \node (check2) [decision, below of=tier2] {Conf > 0.75?};
    \node (tier3) [tier, below of=check2] {Tier 3: LLM};
    \node (ensemble) [tier, left of=tier3, xshift=-3.5cm] {Ensemble Voting};
    \node (output) [result, below of=ensemble] {Final Classification};
    
    \draw [->, thick] (input) -- (tier1);
    \draw [->, thick] (tier1) -- (check1);
    \draw [->, thick] (check1) -- node[anchor=south, font=\small] {No} (tier2);
    \draw [->, thick] (check1) -- node[anchor=east, font=\small] {Yes} (ensemble);
    \draw [->, thick] (tier2) -- (check2);
    \draw [->, thick] (check2) -- node[anchor=east, font=\small] {No} (tier3);
    \draw [->, thick] (check2) -- node[anchor=south, font=\small] {Yes} (ensemble);
    \draw [->, thick] (tier3) -- (ensemble);
    \draw [->, thick] (ensemble) -- (output);
\end{tikzpicture}
\caption{3-Tier Cascade Ensemble Architecture}
\label{fig:cascade-ensemble}
\end{figure}

\textbf{Design Rationale:}
\begin{itemize}
    \item \textbf{Rule-Based (Tier 1)}: Fast (50-100ms), low cost, high precision for clear cases
    \item \textbf{BERT (Tier 2)}: Moderate speed (800ms), zero cost, good for ambiguous cases
    \item \textbf{LLM (Tier 3)}: Slow (2-4s), costly (\$0.0001), highest accuracy for complex cases
    \item \textbf{Cascade Logic}: Skip expensive tiers if confidence is high
    \item \textbf{Ensemble Voting}: Combine all invoked models for final decision
\end{itemize}

\section{Tier 1: Rule-Based Classifier}

\subsection{Algorithm Overview}

Keyword-based scoring system with MeSH term matching and structural analysis.

\textbf{Classification Logic:}
\begin{lstlisting}[style=pythonstyle]
# services/pv-classification-service/rule_based_classifier.py
class RuleBasedClassifier:
    def __init__(self):
        # High-value ICSR keywords (weight: 5.0)
        self.icsr_keywords = {
            'adverse event': 5.0,
            'adverse reaction': 5.0,
            'side effect': 4.5,
            'patient reported': 4.0,
            'case report': 4.0,
            'case series': 3.5,
            'individual case': 5.0,
            'toxicity': 3.0,
            'anaphylaxis': 4.5,
            'drug-induced': 4.0
        }
        
        # Safety information keywords (weight: 3.0)
        self.safety_keywords = {
            'clinical trial': 3.0,
            'safety analysis': 3.5,
            'post-marketing': 3.5,
            'surveillance': 3.0,
            'aggregate data': 2.5,
            'safety profile': 3.0,
            'risk assessment': 3.0
        }
        
        # MeSH term weights (weight: 3.0)
        self.mesh_weights = {
            'Drug-Related Side Effects and Adverse Reactions': 5.0,
            'Adverse Drug Reaction Reporting Systems': 4.5,
            'Pharmacovigilance': 3.5,
            'Product Surveillance, Postmarketing': 3.0
        }
        
        # Publication type weights (weight: 4.0)
        self.pub_type_weights = {
            'Case Reports': 4.5,
            'Clinical Trial': 3.0,
            'Observational Study': 2.5
        }
    
    def classify(self, document: dict) -> dict:
        """Classify document using rule-based scoring"""
        score = 0.0
        evidence = []
        
        # 1. Keyword matching in title/abstract
        text = (document.get('title', '') + ' ' + 
                document.get('abstract', '')).lower()
        
        for keyword, weight in self.icsr_keywords.items():
            if keyword in text:
                score += weight
                evidence.append(f"ICSR keyword: '{keyword}' (+{weight})")
        
        for keyword, weight in self.safety_keywords.items():
            if keyword in text:
                score += weight * 0.6  # Lower weight for safety keywords
                evidence.append(f"Safety keyword: '{keyword}' (+{weight*0.6})")
        
        # 2. MeSH term matching
        mesh_terms = document.get('mesh_terms', [])
        for mesh in mesh_terms:
            if mesh in self.mesh_weights:
                weight = self.mesh_weights[mesh]
                score += weight
                evidence.append(f"MeSH term: '{mesh}' (+{weight})")
        
        # 3. Publication type matching
        pub_types = document.get('publication_types', [])
        for pub_type in pub_types:
            if pub_type in self.pub_type_weights:
                weight = self.pub_type_weights[pub_type]
                score += weight
                evidence.append(f"Pub type: '{pub_type}' (+{weight})")
        
        # 4. Structural analysis
        if self._has_patient_section(document):
            score += 2.0
            evidence.append("Has patient section (+2.0)")
        
        if self._has_temporal_markers(document):
            score += 1.5
            evidence.append("Has temporal markers (+1.5)")
        
        if self._has_causality_assessment(document):
            score += 2.5
            evidence.append("Has causality assessment (+2.5)")
        
        # 5. Classification decision
        if score >= 12.0:
            classification = "ICSR"
            confidence = min(0.95, 0.7 + (score - 12) * 0.02)
        elif score >= 7.0:
            classification = "Safety_Information"
            confidence = 0.7 + (score - 7) * 0.03
        elif score >= 3.0:
            classification = "Non-ICSR"
            confidence = 0.6 + (score - 3) * 0.025
        else:
            classification = "Rejected"
            confidence = 0.5 + score * 0.05
        
        return {
            'classification': classification,
            'confidence': confidence,
            'score': score,
            'evidence': evidence,
            'processing_time_ms': 75  # Avg processing time
        }
\end{lstlisting}

\subsection{Performance Metrics}

\begin{table}[h]
\centering
\caption{Rule-Based Classifier Performance}
\begin{tabularx}{\textwidth}{lXXXX}
\toprule
\textbf{Category} & \textbf{Precision} & \textbf{Recall} & \textbf{F1 Score} & \textbf{Support} \\
\midrule
ICSR & 0.89 & 0.76 & 0.82 & 120 \\
Safety\_Information & 0.75 & 0.82 & 0.78 & 180 \\
Non-ICSR & 0.81 & 0.87 & 0.84 & 150 \\
Rejected & 0.92 & 0.84 & 0.88 & 50 \\
\textbf{Overall} & \textbf{0.82} & \textbf{0.82} & \textbf{0.82} & \textbf{500} \\
\bottomrule
\end{tabularx}
\end{table}

\textbf{Strengths:}
\begin{itemize}
    \item Extremely fast (50-100ms per document)
    \item Zero API cost
    \item Transparent (explainable evidence)
    \item High precision for keyword-rich documents
\end{itemize}

\textbf{Weaknesses:}
\begin{itemize}
    \item Struggles with novel terminology
    \item Low recall for nuanced cases
    \item Cannot understand context
    \item Fixed rules (not adaptive)
\end{itemize}

\section{Tier 2: BERT Zero-Shot Classifier}

\subsection{Model Architecture}

Uses pre-trained `facebook/bart-large-mnli` for natural language inference.

\textbf{Implementation:}
\begin{lstlisting}[style=pythonstyle]
# services/pv-classification-service/bert_classifier.py
from transformers import pipeline
import torch

class BERTClassifier:
    def __init__(self):
        # Load pre-trained zero-shot model
        self.model = pipeline(
            "zero-shot-classification",
            model="facebook/bart-large-mnli",
            device=0 if torch.cuda.is_available() else -1  # Use GPU if available
        )
        
        # Classification hypotheses
        self.hypotheses = [
            "This is an individual case safety report describing adverse events in a specific patient.",
            "This is general drug safety information or aggregate adverse event data.",
            "This is medical research without focus on drug safety or adverse events.",
            "This is not relevant to pharmacovigilance or drug safety."
        ]
        
        self.class_mapping = {
            0: "ICSR",
            1: "Safety_Information",
            2: "Non-ICSR",
            3: "Rejected"
        }
    
    def classify(self, document: dict) -> dict:
        """Classify using BERT zero-shot"""
        # Prepare input text (title + abstract, max 512 tokens)
        text = document.get('title', '') + '. ' + document.get('abstract', '')
        text = text[:2000]  # Truncate to reasonable length
        
        # Run zero-shot classification
        result = self.model(
            text,
            candidate_labels=self.hypotheses,
            multi_label=False  # Single label classification
        )
        
        # Extract top prediction
        top_idx = 0
        top_score = result['scores'][0]
        classification = self.class_mapping[top_idx]
        
        # Adjust confidence (BERT tends to be overconfident)
        confidence = top_score * 0.9  # 10% penalty
        
        return {
            'classification': classification,
            'confidence': confidence,
            'scores': {
                'ICSR': result['scores'][0],
                'Safety_Information': result['scores'][1],
                'Non-ICSR': result['scores'][2],
                'Rejected': result['scores'][3]
            },
            'processing_time_ms': 850  # Avg processing time
        }
\end{lstlisting}

\subsection{Performance Metrics}

\begin{table}[h]
\centering
\caption{BERT Classifier Performance}
\begin{tabularx}{\textwidth}{lXXXX}
\toprule
\textbf{Category} & \textbf{Precision} & \textbf{Recall} & \textbf{F1 Score} & \textbf{Support} \\
\midrule
ICSR & 0.76 & 0.68 & 0.72 & 120 \\
Safety\_Information & 0.71 & 0.79 & 0.75 & 180 \\
Non-ICSR & 0.82 & 0.75 & 0.78 & 150 \\
Rejected & 0.85 & 0.70 & 0.77 & 50 \\
\textbf{Overall} & \textbf{0.78} & \textbf{0.73} & \textbf{0.75} & \textbf{500} \\
\bottomrule
\end{tabularx}
\end{table}

\textbf{Strengths:}
\begin{itemize}
    \item Understands semantic meaning
    \item Handles novel terminology better than rules
    \item Zero API cost (local model)
    \item Reasonable speed with GPU
\end{itemize}

\textbf{Weaknesses:}
\begin{itemize}
    \item Lower precision than rules
    \item Overconfident predictions (requires calibration)
    \item Black box (not explainable)
    \item Requires significant memory (1.5 GB model)
\end{itemize}

\section{Tier 3: LLM Classifier (GPT-4o-mini)}

\subsection{Expert Prompt Engineering}

150-line carefully crafted prompt with examples and guidelines.

\textbf{Prompt Structure:}
\begin{lstlisting}[style=pythonstyle]
PHARMACOVIGILANCE_EXPERT_PROMPT = """
You are a senior pharmacovigilance expert with 20+ years of experience in adverse event reporting and drug safety.

CRITICAL CLASSIFICATION CRITERIA:

1. ICSR (Individual Case Safety Report):
   MUST contain ALL of the following:
   - Specific identifiable patient (even if anonymized: "55-year-old female")
   - Drug/vaccine/medical product administered
   - Adverse event with temporal relationship
   - Clinical outcome (recovered, ongoing, fatal, etc.)
   
   Example ICSR:
   "A 45-year-old female with no prior cardiac history developed Stevens-Johnson 
   syndrome 3 days after starting allopurinol 300mg daily for gout. The patient 
   was hospitalized and recovered after drug discontinuation and corticosteroid 
   treatment."

2. Safety_Information:
   Contains safety data WITHOUT individual patient details:
   - Aggregate adverse event rates
   - Clinical trial safety results
   - Post-marketing surveillance summaries
   - Drug label warnings
   
   Example Safety_Information:
   "In a phase III trial of 1,200 patients, Drug X showed a 15% incidence of 
   nausea and 8% incidence of headache. Serious adverse events occurred in 
   2% of patients."

3. Non-ICSR:
   Medical content without safety focus:
   - Efficacy studies
   - Mechanism of action research
   - Epidemiological studies (non-safety)
   - Prescribing guidelines
   
   Example Non-ICSR:
   "Drug X inhibits COX-2 enzyme selectively, providing anti-inflammatory 
   effects. Efficacy was demonstrated in 80% of patients with arthritis."

4. Rejected:
   Not relevant to pharmacovigilance:
   - Financial reports
   - Opinion pieces
   - Non-medical topics
   
CONFIDENCE SCORING GUIDELINES:
- 0.95-1.00: Textbook example, all criteria clearly met
- 0.85-0.94: Strong evidence, minor ambiguity
- 0.75-0.84: Moderate evidence, some missing details
- 0.65-0.74: Weak evidence, significant uncertainty
- <0.65: Very uncertain, consider alternative classification

OUTPUT FORMAT (JSON):
{
  "classification": "ICSR" | "Safety_Information" | "Non-ICSR" | "Rejected",
  "confidence": 0.0-1.0,
  "reasoning": "Detailed 2-3 sentence explanation",
  "key_entities": ["patient_descriptor", "drug_name", "adverse_event", "outcome"]
}

Now classify the following document:
Title: {title}
Abstract: {abstract}
"""
\end{lstlisting}

\subsection{Implementation}

\begin{lstlisting}[style=pythonstyle]
# services/pv-classification-service/llm_classifier.py
from openai import OpenAI
import json

class LLMClassifier:
    def __init__(self):
        self.client = OpenAI(api_key=os.getenv("OPENAI_API_KEY"))
        self.model = "gpt-4o-mini"
        self.temperature = 0.5  # Balanced creativity/consistency
    
    def classify(self, document: dict) -> dict:
        """Classify using GPT-4o-mini with expert prompt"""
        # Prepare prompt
        prompt = PHARMACOVIGILANCE_EXPERT_PROMPT.format(
            title=document.get('title', ''),
            abstract=document.get('abstract', '')
        )
        
        # Call OpenAI API
        response = self.client.chat.completions.create(
            model=self.model,
            messages=[
                {"role": "system", "content": "You are a pharmacovigilance expert."},
                {"role": "user", "content": prompt}
            ],
            response_format={"type": "json_object"},  # Force JSON output
            temperature=self.temperature,
            max_tokens=500
        )
        
        # Parse response
        result_json = json.loads(response.choices[0].message.content)
        
        # Calculate cost
        input_tokens = response.usage.prompt_tokens
        output_tokens = response.usage.completion_tokens
        cost = (input_tokens * 0.00015 + output_tokens * 0.0006) / 1000  # GPT-4o-mini pricing
        
        return {
            'classification': result_json['classification'],
            'confidence': result_json['confidence'],
            'reasoning': result_json['reasoning'],
            'entities': result_json.get('key_entities', []),
            'processing_time_ms': 2500,  # Avg API latency
            'cost': cost,
            'tokens_used': input_tokens + output_tokens
        }
\end{lstlisting}

\subsection{Performance Metrics}

\begin{table}[h]
\centering
\caption{LLM Classifier Performance}
\begin{tabularx}{\textwidth}{lXXXX}
\toprule
\textbf{Category} & \textbf{Precision} & \textbf{Recall} & \textbf{F1 Score} & \textbf{Support} \\
\midrule
ICSR & 0.94 & 0.90 & 0.92 & 120 \\
Safety\_Information & 0.88 & 0.91 & 0.89 & 180 \\
Non-ICSR & 0.91 & 0.89 & 0.90 & 150 \\
Rejected & 0.96 & 0.88 & 0.92 & 50 \\
\textbf{Overall} & \textbf{0.92} & \textbf{0.90} & \textbf{0.91} & \textbf{500} \\
\bottomrule
\end{tabularx}
\end{table}

\textbf{Strengths:}
\begin{itemize}
    \item Highest accuracy across all categories
    \item Handles complex, nuanced cases
    \item Provides detailed reasoning
    \item Extracts key entities
    \item Adaptable via prompt engineering
\end{itemize}

\textbf{Weaknesses:}
\begin{itemize}
    \item Slow (2-4 seconds per document)
    \item Costly (\$0.0001 per classification)
    \item Requires API key
    \item Non-deterministic (temperature > 0)
    \item Potential for hallucinations
\end{itemize}

\section{Ensemble Voting Mechanism}

\subsection{Weighted Voting Algorithm}

Combines predictions from all invoked models using weighted voting.

\begin{lstlisting}[style=pythonstyle]
# services/pv-classification-service/ensemble_classifier.py
class EnsembleClassifier:
    def __init__(self, rule_weight=0.48, bert_weight=0.30, llm_weight=0.27):
        self.rule_classifier = RuleBasedClassifier()
        self.bert_classifier = BERTClassifier()
        self.llm_classifier = LLMClassifier()
        
        # Weights learned via RL optimization
        self.weights = {
            'rule_based': rule_weight,
            'bert': bert_weight,
            'llm': llm_weight
        }
        
        # Confidence thresholds for cascade
        self.rule_threshold = 0.85
        self.bert_threshold = 0.75
    
    def classify(self, document: dict, force_ensemble: bool = False) -> dict:
        """Classify document using cascade ensemble"""
        results = {}
        decision_path = []
        
        # Tier 1: Rule-based (always run)
        rule_result = self.rule_classifier.classify(document)
        results['rule_based'] = rule_result
        decision_path.append('rule_based')
        
        # Check if we can skip BERT/LLM
        if not force_ensemble and rule_result['confidence'] >= self.rule_threshold:
            # High confidence from rules, skip other models
            return self._finalize_result(results, decision_path, document)
        
        # Tier 2: BERT (if rule confidence is low)
        bert_result = self.bert_classifier.classify(document)
        results['bert'] = bert_result
        decision_path.append('bert')
        
        # Check if we can skip LLM
        if not force_ensemble and bert_result['confidence'] >= self.bert_threshold:
            # Sufficient confidence from BERT, skip LLM
            return self._finalize_result(results, decision_path, document)
        
        # Tier 3: LLM (for complex cases)
        llm_result = self.llm_classifier.classify(document)
        results['llm'] = llm_result
        decision_path.append('llm')
        
        return self._finalize_result(results, decision_path, document)
    
    def _finalize_result(self, results: dict, decision_path: list, document: dict) -> dict:
        """Compute weighted ensemble vote"""
        # Weighted scores for each class
        vote_scores = defaultdict(float)
        
        for model_name, result in results.items():
            classification = result['classification']
            confidence = result['confidence']
            weight = self.weights[model_name]
            
            vote_scores[classification] += weight * confidence
        
        # Normalize to probabilities
        total = sum(vote_scores.values())
        probabilities = {k: v/total for k, v in vote_scores.items()}
        
        # Final classification
        final_class = max(probabilities, key=probabilities.get)
        final_confidence = probabilities[final_class]
        
        # Calculate total cost and time
        total_cost = sum(r.get('cost', 0) for r in results.values())
        total_time = sum(r.get('processing_time_ms', 0) for r in results.values())
        
        return {
            'classification': final_class,
            'classification_confidence': final_confidence,
            'classification_reasoning': results.get('llm', {}).get('reasoning', 'Ensemble decision'),
            'classification_entities': results.get('llm', {}).get('entities', []),
            'ensemble_metadata': {
                'models_used': decision_path,
                'individual_results': results,
                'decision_path': decision_path,
                'processing_time_ms': total_time,
                'classification_cost': total_cost
            }
        }
\end{lstlisting}

\subsection{Ensemble Performance}

\begin{table}[h]
\centering
\caption{Model Comparison - Individual vs. Ensemble}
\begin{tabularx}{\textwidth}{lXXXXX}
\toprule
\textbf{Model} & \textbf{Precision} & \textbf{Recall} & \textbf{F1 Score} & \textbf{Cost/1K} & \textbf{Time/doc} \\
\midrule
Rule-Based & 0.82 & 0.82 & 0.82 & \$0.00 & 75ms \\
BERT & 0.78 & 0.73 & 0.75 & \$0.00 & 850ms \\
LLM & 0.92 & 0.90 & 0.91 & \$100 & 2500ms \\
\textbf{Ensemble (RL-opt)} & \textbf{0.91} & \textbf{0.89} & \textbf{0.90} & \textbf{\$2.95} & \textbf{1800ms} \\
\bottomrule
\end{tabularx}
\end{table}

\textbf{Key Findings:}
\begin{itemize}
    \item Ensemble achieves 99\% of LLM accuracy at 3\% of the cost
    \item 28\% faster than LLM-only approach
    \item Cascade logic skips LLM for 65\% of documents
    \item RL optimization boosted rule weight (+60\%), penalized BERT (-38\%)
\end{itemize}

\begin{figure}[h]
\centering
\fbox{\textit{[SCREENSHOT PLACEHOLDER: Model Performance Comparison - Bar Chart]}}
\caption{Precision/Recall/F1 Comparison Across Models}
\label{fig:model-performance-comparison}
\end{figure}

\section{Classification Categories}

\subsection{ICSR (Individual Case Safety Report)}

\textbf{Definition:}
A detailed report of a suspected adverse reaction to a drug in a specific, identifiable patient.

\textbf{Required Elements:}
\begin{enumerate}
    \item \textbf{Patient}: Age, gender, medical history (can be anonymized)
    \item \textbf{Drug}: Name, dosage, route, start/end dates
    \item \textbf{Adverse Event}: Description, onset, severity
    \item \textbf{Temporal Relationship}: Time between drug and event
    \item \textbf{Outcome}: Recovered, ongoing, fatal, hospitalized
\end{enumerate}

\textbf{Examples:}
\begin{itemize}
    \item "A 67-year-old male developed acute liver failure 10 days after starting atorvastatin..."
    \item "Patient presented with anaphylaxis 30 minutes after first dose of penicillin..."
\end{itemize}

\subsection{Safety\_Information}

\textbf{Definition:}
General drug safety information without individual patient cases.

\textbf{Includes:}
\begin{itemize}
    \item Aggregate adverse event statistics
    \item Clinical trial safety results
    \item Post-marketing surveillance data
    \item Drug label safety sections
    \item Pharmacoepidemiology studies
\end{itemize}

\textbf{Examples:}
\begin{itemize}
    \item "In 1,200 patients treated with Drug X, 15\% experienced nausea..."
    \item "Post-marketing surveillance revealed 50 cases of hepatotoxicity over 5 years..."
\end{itemize}

\subsection{Non-ICSR}

\textbf{Definition:}
Medical literature without drug safety focus.

\textbf{Includes:}
\begin{itemize}
    \item Efficacy studies
    \item Mechanism of action research
    \item Pharmacokinetics/pharmacodynamics
    \item Drug discovery
    \item Treatment guidelines
\end{itemize}

\textbf{Examples:}
\begin{itemize}
    \item "Drug X showed 80\% response rate in treating hypertension..."
    \item "Molecular docking studies revealed Drug X binds to enzyme Y..."
\end{itemize}

\subsection{Rejected}

\textbf{Definition:}
Not relevant to pharmacovigilance.

\textbf{Includes:}
\begin{itemize}
    \item Financial reports
    \item Opinion pieces / editorials
    \item Non-medical topics
    \item Duplicate publications
\end{itemize}

\section{Model Retraining Pipeline}

\subsection{Data Collection for Retraining}

\begin{commandbox}
\begin{lstlisting}[style=bashstyle]
# Collect user-corrected classifications
python ml_training/collect_training_data.py \
    --source elasticsearch \
    --index documents \
    --filter user_corrected=true \
    --min-samples 1000 \
    --output ml_training/data/retraining_data.csv
\end{lstlisting}
\end{commandbox}

\subsection{Fine-Tuning BERT}

\begin{lstlisting}[style=pythonstyle]
# ml_training/train_bert.py
from transformers import Trainer, TrainingArguments

# Load collected data
df = pd.read_csv('data/retraining_data.csv')

# Fine-tune BERT model
training_args = TrainingArguments(
    output_dir='./bert_finetuned',
    num_train_epochs=3,
    per_device_train_batch_size=8,
    evaluation_strategy="steps",
    save_steps=500,
    learning_rate=2e-5
)

trainer = Trainer(
    model=model,
    args=training_args,
    train_dataset=train_dataset,
    eval_dataset=eval_dataset
)

trainer.train()
\end{lstlisting}

\subsection{Model Evaluation \& A/B Testing}

\begin{lstlisting}[style=pythonstyle]
# Evaluate new model on held-out test set
python ml_training/evaluate_model.py \
    --model ml_training/bert_finetuned \
    --test-data ml_training/data/test_set.csv \
    --output ml_training/results/evaluation_report.json

# A/B test: Route 10% of traffic to new model
python ml_training/ab_test.py \
    --model-a ensemble_v1 \
    --model-b ensemble_v2_with_finetuned_bert \
    --traffic-split 0.9,0.1 \
    --duration 7days
\end{lstlisting}

\subsection{Deployment}

\begin{commandbox}
\begin{lstlisting}[style=bashstyle]
# Deploy new model if A/B test shows 2%+ improvement
docker cp ml_training/bert_finetuned pv-classification-service:/app/models/
docker-compose restart pv-classification-service

# Verify deployment
curl http://localhost:8012/api/v1/model/info
\end{lstlisting}
\end{commandbox}

\end{document}

